\documentclass[11pt,letterpaper]{article}
\usepackage[margin=0.88in,headheight=15pt]{geometry}
\usepackage[T1]{fontenc}
\usepackage{lmodern,microtype,amsmath,amssymb,amsthm,mathtools}
\usepackage{booktabs,tabularx,array,xcolor,fancyhdr}
\usepackage[colorlinks=true,linkcolor=blue!45!black,urlcolor=blue!50!black,citecolor=blue!45!black]{hyperref}
\definecolor{ink}{RGB}{28,49,68}
\pagestyle{fancy}\fancyhf{}
\setlength{\parindent}{0pt}\setlength{\parskip}{5pt}
\setlength{\emergencystretch}{2em}
\setcounter{tocdepth}{1}\numberwithin{equation}{section}
\theoremstyle{plain}
\newtheorem{theorem}{Theorem}[section]
\newtheorem{lemma}[theorem]{Lemma}
\newtheorem{proposition}[theorem]{Proposition}
\newtheorem{corollary}[theorem]{Corollary}
\theoremstyle{remark}\newtheorem{remark}[theorem]{Scope}
\DeclareMathOperator{\Var}{Var}\DeclareMathOperator{\Cov}{Cov}
\DeclareMathOperator{\Tr}{Tr}\DeclareMathOperator{\Vol}{Vol}
\DeclareMathOperator{\supp}{supp}\DeclareMathOperator{\diver}{div}
\DeclareMathOperator{\osc}{osc}\DeclareMathOperator{\Ad}{Ad}
\newcommand{\E}{\mathbb E}\newcommand{\R}{\mathbb R}
\newcommand{\dd}{\,\mathrm d}\newcommand{\OH}{\mathcal O_H}
\newcommand{\norm}[1]{\left\lVert#1\right\rVert}
\newcommand{\abs}[1]{\left|#1\right|}
\newcommand{\ind}{\mathbf1}
\title{Physical Wilson contacts and local conditional subtraction}\author{}\date{}
\begin{document}\maketitle
\section{Conventions, the unchanged source, and its actual contact}
\label{sec:conventions}

Let $G=SU(N)$ with $N\ge2$, $m=\dim G=N^2-1$, with normalized product Haar measure
$\lambda$. Except where an explicit action convention is displayed, use
\begin{equation}
s_p=N-\Re\Tr U_p,\quad S_0=\sum_p s_p,\quad
\rho=e^{-\beta S_0},\quad Z_\beta=\int\rho\dd\lambda,\quad
\dd\mu_\beta=Z_\beta^{-1}\rho\dd\lambda,
\qquad\beta=c_\beta g^{-2}.
\label{eq:Wilson}
\end{equation}
The physical fine spacing is $a$. A finite source algebra has commuting generators
$\eta_i$ with $\eta_i^2=0$; write $\eta_I=\prod_{i\in I}\eta_i$.
All exponentials and logarithms in this algebra are finite coefficientwise operations.

Use the original invariant derivative
$X_{e,A}F(U)=\left.\frac{\dd}{\dd t}F(\ldots,e^{tT_A}U_e,\ldots)\right|_{t=0}$.
This explicit convention avoids ambiguity in naming left and right invariant fields.
Put
\begin{align}
F_i&=X_iS_0=\sum_{\mu\ne\nu_i}F_{i;\mu},\notag\\
D[f]&=\sum_i f_i\sum_{\mu\ne\nu_i}F_{i;\mu}^2,
&P[f]&=\sum_i f_iF_i^2,
&\ell[f]&=\sum_i f_iX_i^2S_0,\notag\\
\sigma&=g^{-2}a^{-2},
&\OH(f)&=\sigma\{D[f]-P[f]/3\}.
\label{eq:H}
\end{align}
Every pointed test offset and original finite hypercubic average is retained.
When scalar formulas are written for a representative, the final expression is
the corresponding original average; degree two uses its double average.

The original bare source is $\exp(\sum_i\eta_i\OH(f_i))$, with no added bare
mixed action coefficient. For two labels, a later actual source action is
\begin{equation}
A=\eta_1O_1+\eta_2O_2+\eta_1\eta_2A_{12}.
\label{eq:source}
\end{equation}
$A_{12}$ must be retained after it has been generated.

For a vector field $V$, define its density adjoint by
$D_Vu=-\diver_\lambda(uV)$. The original PD12 decomposition is
\begin{align}
V_H(f)&=-\frac{\sigma}{3\beta}\sum_i f_iF_iX_i,
&t_H(f)&=-\frac{\sigma}{3}\{P[f]-\beta^{-1}\ell[f]\},\notag\\
r_H(f)&=\sigma\{D[f]-(3\beta)^{-1}\ell[f]\},
&\OH(f)&=t_H(f)+r_H(f).
\label{eq:oldH}
\end{align}
In physical units the ordinary part has the exact rows
\begin{align}
r_H(f)&=c_f+d_f+\ell_f,\notag\\
d_f&=\gamma^{-2}D_a(f),\qquad
\ell_f=\frac{C_F}{3\beta_s}\gamma^{-2}a^{-2}\sum_p f_ps_p,\notag\\
c_f&=-\frac{C_F}{3\beta_s}\gamma^{-2}a^{-2}N\sum_p f_p.
\label{eq:Casimir}
\end{align}
Here $f_p$ is the literal sum of the four original bond incidences, including the
specified shifts. The notation $\gamma,\beta_s$ allows saved source normalizations
to remain distinct from a later native coupling. At the original birth these
specialize to $\gamma=g$, $\beta_s=\beta$.
Here $D_a(f)=a^{-2}D[f]=a^4\sum_{b,A,\mu\ne\nu_b}f_b(J_{b;\mu}^A)^2$
with $J_{b;\mu}=a^{-3}F_{b;\mu}$.

If $O_i=t_i+r_i$ and $t_i\rho=D_{V_i}\rho$, the original compensated coefficient is
\begin{equation}
b_{12}=A_{12}+\tfrac12(V_1t_2+V_2t_1)+V_1r_2+V_2r_1.
\label{eq:oldb}
\end{equation}
For the bare $H$ pair PD12 gives, with $W_f=\sum_i f_iF_iX_i$,
\begin{equation}
b_{HH}(f,h)=\frac{\sigma^2}{9\beta}
\left[\tfrac12(W_fP[h]+W_hP[f])-3(W_fD[h]+W_hD[f])
+\frac{W_f\ell[h]+W_h\ell[f]}{2\beta}\right].
\label{eq:bareb}
\end{equation}
Its last term is the Casimir response. It is not dropped in any estimate below.

\section{Actual Wilson local moments and an improved original contact bound}
\label{sec:moments}

The useful local input is stronger than the crude
$\E s_p\le C\log(2+\beta)/\beta$ used in~\cite[pp.~21--23]{W}. We rederive the
matching partition powers from~\cite[\S10, pp.~21--23]{W}, keeping their finite-size loss.
Let $n$ be the number of fine sites per side, and take an even periodic torus.

\begin{lemma}[Matching finite-volume partition powers]
\label{lem:partition}
There are group- and action-dependent constants $0<c_l\le1\le c_u$ such that
for $\beta\ge1$,
\begin{equation}
(c_l\beta^{-m/2})^{Q_n}\le Z_{n,\beta}
\le(c_u\beta^{-m/2})^{R_n},\quad
Q_n=3n^4+1,\quad R_n=3n^4-4n^3+1.
\label{eq:partition}
\end{equation}
The upper estimate holds at every positive coefficient after enlarging $c_u$.
\end{lemma}
\begin{proof}
Gauge-fixing a spanning tree by Haar-invariant changes of variables leaves exactly
$4n^4-(n^4-1)=Q_n$ Haar links. Restrict them to radius $c\beta^{-1/2}$ about the
identity. Each plaquette then has $\beta s_p\le C$, and each retained link ball
has mass at least $c'\beta^{-m/2}$. The action cost is at most a fixed constant
per retained link, giving the lower bound with its original gauge factor one.

For the upper bound, discard the nonnegative actions crossing the periodic seam.
On the resulting free box, fix a temporal tree and a recursive tree on its bottom
face. Keep the temporal plaquettes and the bottom-face plaquettes. Integrating
each spatial-link chain from the free end gives the exact factor
$z(\beta)=\int_Ge^{-\beta s(U)}\dd U$. The unique nondegenerate action minimum
and compactness give $z(\beta)\le c_u\beta^{-m/2}$. In dimension $d$ the number
of these factors is
$R_d(n)=(d-1)(n-1)^2n^{d-2}+R_{d-1}(n)$, $R_1=0$.
Induction gives $R_d=(d-1)n^d-dn^{d-1}+1$ and hence the stated $R_n$.
No discarded action or gauge integral contributes an unrecorded inverse mass.
\end{proof}

\begin{proposition}[Local moments with their finite-size factor]
\label{prop:moments}
For every even $n$, $\beta\ge1$, and original plaquette $p$, there is a fixed
$C_G$ for which
\begin{align}
\E_{n,\beta}e^{\beta s_p/8}&\le\mathcal B_{n,\beta},\notag\\
\mathcal B_{n,\beta}
&=c_u^{R_n/n^4}c_l^{-Q_n/n^4}
2^{mR_n/(2n^4)}\beta^{2m/n}
\le C_G\exp\!\left(\frac{2m\log\beta}{n}\right),\label{eq:B}\\
\E_{n,\beta}(\beta s_p)^r&\le r!8^r\mathcal B_{n,\beta}
\quad(r\ge1).
\label{eq:moment}
\end{align}
In particular this factor is uniformly bounded when $\log\beta/n$ is bounded.
\end{proposition}
\begin{proof}
Use original site-plane Wilson reflection positivity and its chessboard inequality
on closed unit cells. The latter is the standard repeated-reflection estimate
for functions supported on a closed block; even cell counts suffice~\cite{biskup}.
It applies to the bounded gauge-invariant function $e^{t\beta s_p}$ on a face of
one unit cell. Its disseminated plaquettes have multiplicity at most four: the
plaquette occupies the full cell in its two directions and can be repeated twice
in each transverse direction. Thus their summed action $S_*\le4S_0$.
There are $n^4$ cells, so for $0\le t<1/4$,
\[
\E e^{t\beta s_p}
\le\left[\E\exp(t\beta S_*)\right]^{1/n^4}
\le\left[\frac{Z_{n,(1-4t)\beta}}{Z_{n,\beta}}\right]^{1/n^4}.
\]
Take $t=1/8$ and apply Lemma~\ref{lem:partition}; the power of $\beta$ is
$m(Q_n-R_n)/(2n^4)=2m/n$. The moment bound follows from
$x^r\le r!8^r e^{x/8}$. Site-plane reflection positivity here is for the original
Wilson measure: conditioned on the common plane links the two reflected halves
factor with conjugate weights. No history-conditioned law is assumed reflection
positive. Alternatively, on tori divisible by eight, a side-four cell with the
plaquette strictly inside gives the original \S10 estimate at $t=1/2$, with its
explicit factor $\beta^{2m4^4/n}$.
\end{proof}

If $Y$ is any literal finite plaquette set, positivity and the elementary power-mean
inequality give
\begin{equation}
\E\left(\beta\sum_{p\in Y}s_p\right)^r
\le |Y|^r r!8^r\mathcal B_{n,\beta}.
\label{eq:patchmoment}
\end{equation}
No independence of the plaquettes is asserted. In the simultaneous regime
$n=t_*/a$, $\beta=O(\log(1/a))$, $\mathcal B_{n,\beta}$ is bounded. This is
not a uniform-in-$\beta$ assertion at fixed $n$, and it does not construct the
renormalized coupling trajectory.

\begin{corollary}[Improved absolute mass of the unchanged bare $H$ contact]
\label{cor:originalcontact}
Let $Y=Y(f,h)$ be the original overlap union in PD12.26 and $N_Y=|Y|$.
For the original source and Wilson measure,
\begin{equation}
\E|b_{HH}(f,h)|
\le C_Gc_\beta^{-2}a^{-4}N_Y
\|f\|_\infty\|h\|_\infty\mathcal B_{n,\beta},\qquad\beta\ge1.
\label{eq:originalcontactmass}
\end{equation}
The corresponding raw mass is this bound multiplied by the actual $Z_\beta$.
\end{corollary}
\begin{proof}
The exact compact-group estimate PD12.27, including every term of
\eqref{eq:bareb}, is
$|b_{HH}|\le C_Gc_\beta^{-2}a^{-4}(1+\beta^{-1})\|f\|_\infty\|h\|_\infty E_Y$,
where $E_Y=\beta\sum_{p\in Y}s_p$. Apply \eqref{eq:patchmoment} with $r=1$.
The one-reserve estimate itself follows from
$|X_is_p|\le C_G\sqrt{s_p}$, bounded fixed second plaquette derivatives,
and the original finite star incidence, as in PD12.28.
\end{proof}
This improves a real birth-contact estimate for the unchanged physical source.
It is not the all-history bound for the ordinary covariance accompanying that
contact. In particular the Casimir action row in \eqref{eq:Casimir} remains.

\begin{corollary}[Sharper error in the nonlinear dispersion theorem]
\label{cor:dispersion}
On the compatible even-torus family, the original nonlinear source of~\cite[\S2]{ND}
satisfies, with its fixed plateau test and local patch-fit assumptions, for
$\beta\ge\beta_0$, $0<a<a_0$, and $L\ge0$,
\begin{equation}
\sup_c\mu_\beta\{|\OH(f)-c|\le L\}
\le\min\{1,C_f[e^{-c_f'\beta}+a^2+a^4L]\}.
\label{eq:sharperdispersion}
\end{equation}
\end{corollary}
\begin{proof}
The bath event of~\cite[\S2.1]{ND} is implied by $s_p\le\epsilon^2$ on a fixed finite
patch, with $\epsilon>0$ fixed. Exponential Markov and \eqref{eq:B} give a
complement probability at most $C\mathcal B_{n,\beta}e^{-\beta\epsilon^2/8}$.
Since $n\ge2$, $\mathcal B_{n,\beta}\le C\beta^m$, and the polynomial is absorbed
into $C_\epsilon e^{-\beta\epsilon^2/16}$. Substitute this error in the identical
good-bath count and conditional concentration proof of~\cite[\S\S2.1--2.4]{ND}. All remaining
terms and the original source normalization are unchanged.
\end{proof}
This is a fixed-patch estimate. A union over $O(a^{-4})$ patches need not be small
on a logarithmic coupling trajectory; no global small-field event is inferred.

\section{An exact local conditional contact prescription}
\label{sec:conditional}

Choose a deterministic sparse link set $\mathcal F$ such that no original
plaquette and no complete determining stencil of an $H$ summand reads two
selected links. Fixed finite-range packing gives such a set with positive lattice
density. All original pointed averages are included when choosing the separation.
Let $b$ be the remaining links. Then the exact Wilson measure has
\begin{align}
\rho\dd\lambda
&=e^{-\beta S_{\mathrm{rest}}(b)}\dd b
\prod_{e\in\mathcal F}e^{-\beta S_e(U_e;b)}\dd U_e,\notag\\
d_e(b)&=\int_Ge^{-\beta S_e(U;b)}\dd U,
&K_e(\dd U\mid b)&=d_e(b)^{-1}e^{-\beta S_e(U;b)}\dd U,\notag\\
\rho_b(b)&=e^{-\beta S_{\mathrm{rest}}(b)}\prod_e d_e(b),
&Z_\beta&=\int\rho_b(b)\dd b.
\label{eq:disintegration}
\end{align}
Here $S_e$ is the sum of original plaquette actions containing $e$. All conditional
denominators and the absolute partition function are explicit. This is an original
Wilson disintegration; it is not identified with an unverified $R$-operation kernel.

Literal stencil separation gives, for each chosen local source label,
\begin{equation}
O_i=O_{0i}(b)+\sum_{e\in\mathcal F}h_{ei}(U_e;b),\quad
m_i(b)=K_{\mathcal F}O_i,\quad
X_{ei}=h_{ei}-K_eh_{ei},\quad K_{\mathcal F}=\bigotimes_eK_e.
\label{eq:localdecomp}
\end{equation}
Assigning each whole original summand that reads $e$ to $h_{ei}$ is one valid
choice. Bath-only additive changes of the splitting do not change its contacts.

\begin{theorem}[Finite-cutoff conditional normal ordering]
\label{thm:normalorder}
The input here is the original linear bare source: every inherited action
coefficient $A_I$ with $|I|\ge2$ is zero.
Define the bath-dependent polynomial of source degree at least two
\begin{equation}
Q_{\mathcal F}(\eta,b)
=\sum_e\log K_e\exp\left(\sum_i\eta_iX_{ei}\right).
\label{eq:QF}
\end{equation}
It has no linear term. The explicitly \emph{new} source prescription
\begin{equation}
\mathcal N_{\mathcal F}(\eta,U)
=\exp\left(\sum_i\eta_iO_i(U)-Q_{\mathcal F}(\eta,b)\right)
\label{eq:newsource}
\end{equation}
satisfies the exact identity
\begin{equation}
K_{\mathcal F}\mathcal N_{\mathcal F}(\eta)
=\exp\left(\sum_i\eta_i m_i(b)\right).
\label{eq:CNidentity}
\end{equation}
For $|I|\ge2$, its subtraction coefficient is
\begin{equation}
Q_{\mathcal F,I}(b)=\sum_e\kappa_{K_e}(h_{ei}:i\in I),
\label{eq:cumulants}
\end{equation}
and is supported in test space where all labels meet one selected-link determining
patch. These patches have diameter at most $C_{\rm st}a$.
\end{theorem}
\begin{proof}
Conditional independence gives
$\log K_{\mathcal F}e^{\sum\eta_iO_i}=\sum\eta_i m_i+Q_{\mathcal F}$.
The factor $e^{-Q_{\mathcal F}}$ is bath measurable, proving
\eqref{eq:CNidentity}. Expanding the finite logarithm gives the joint conditional
cumulants without a factorial for distinct square-zero labels. A cumulant
vanishes if one of its marks does not read that selected patch, proving locality.
\end{proof}

The original source remains
$K_{\mathcal F}e^{\sum\eta O}=e^{\sum\eta m+Q_{\mathcal F}}$.
Removing $Q_{\mathcal F}$ from this identity changes the source and is legitimate
only under the explicit new definition \eqref{eq:newsource}. Its source-zero mass
is still $Z_\beta$, and its linear insertion is still the original $O_i$.
For an inherited nonlinear source action, its complete $K_{\mathcal F}e^A$ must
instead be used. The one-link cumulant formula \eqref{eq:cumulants} alone would
then omit its inherited coefficients.

At degree two, write
\begin{align}
C_{\mathcal F}(f,h;b)&=\sum_e\Cov_{K_e}(h_e(f),h_e(h)),\notag\\
[O_fO_h]_{\mathcal F}&=O_fO_h-C_{\mathcal F}(f,h;b),\notag\\
\Cov_{\mu_\beta}(O_f,O_h)
&=\Cov_{\mu_\beta}(m_f,m_h)+\E C_{\mathcal F}(f,h;b).
\label{eq:variance-decomp}
\end{align}
The conditional contact is positive semidefinite in its real test labels for every
bath. Its coefficients can depend on neighboring bath links. It is not replaced
by a scalar identity counterterm. The full new hierarchy equals the moment
hierarchy of the projected finite-cutoff observables $m_i$.

Each projected local density is gauge invariant and supported in an enlarged
microscopic stencil; a smeared $m_f$ is supported in the corresponding enlargement
of $\supp f$. Gauge covariance follows by changing the integrated link and its
endpoint frames in \eqref{eq:disintegration}. An arbitrary fine carrier must remain
inside the conditional integral in \eqref{eq:CNidentity}; only bath-measurable
carriers factor out. This limitation is relevant to retained readouts.

\begin{theorem}[All finite-order absolute contact masses]
\label{thm:contactmass}
On the original even torus, let $r=|I|\ge2$ for the bare linear source of
Theorem~\ref{thm:normalorder}, and let $N_{\mathcal F}(I)$ count the selected patches meeting
every original support in $I$. For the original $H$ normalizations on the same
original lattice,
\begin{equation}
\E|Q_{\mathcal F,I}|
\le C_{G,r}c_\beta^{-r}\mathcal B_{n,\beta}a^{-2r}
N_{\mathcal F}(I)\prod_{i\in I}\|f_i\|_\infty.
\label{eq:contactmass}
\end{equation}
The raw absolute mass is bounded by the right side times $Z_\beta$.
For a fixed real $f$ equal to one on a physical box of positive volume, the selected
links can be chosen so that
\begin{equation}
c_fa^{-8}\le \E C_{\mathcal F}(f,f;b)
\le C_f\mathcal B_{n,\beta}a^{-8}
\label{eq:contacttwosided}
\end{equation}
for $a$ sufficiently small and $\beta\ge\beta_0$.
\end{theorem}
\begin{proof}
Compact plaquette estimates give, for the whole-summand choice in
\eqref{eq:localdecomp},
$|h_{ei}|\le C_G\sigma\|f_i\|_\infty S_{Y_e}$,
where $Y_e$ is a fixed-size original plaquette patch. The cumulant partition
formula and conditional H\"older give
\[
|\kappa_{K_e}(h_{ei}:i\in I)|
\le C_{G,r}\sigma^r\prod_i\|f_i\|_\infty\,K_e(S_{Y_e}^r).
\]
Indeed for a partition $\pi$ of $I$, the product
$\prod_{B\in\pi}K_eS_{Y_e}^{|B|}$ is at most $K_eS_{Y_e}^r$;
the sum of its finite cumulant coefficients is absorbed in $C_{G,r}$.
Integrate over the actual bath law, apply \eqref{eq:patchmoment}, and use
$\sigma/\beta=c_\beta^{-1}a^{-2}$. Only the literal common-overlap patches
contribute. This proves \eqref{eq:contactmass}.

For the lower bound, use selected links with their entire determining stencils
inside the plateau. Section~2 of~\cite{ND} proves for the complete dimensionless local
$H$ summand that $\Var_{K_e}(G_e)\ge c\beta^{-2}$ on each good bath. Each selected
link has good-bath probability at least one half, and there are at least
$c_fa^{-4}$ such links. Multiplication by
$\sigma^2=\beta^2a^{-4}/c_\beta^2$ gives at least $c_fa^{-8}$.
This uses no independence of good-bath events. The upper bound is
\eqref{eq:contactmass} with $r=2$ and the actual count $O(a^{-4})$.
\end{proof}

Saved scalar factors multiply \eqref{eq:contactmass} by their literal product of
absolute values. Different saved meshes or pullbacks are not covered merely by
changing a displayed power of $a$: their actual common determining decomposition
and moment bounds must first be verified.

\subsection{Tests of the prescription and its limitations}

A precise counterexample prevents promoting one conditional pass to complete
renormalization. Let independent Gaussian pairs $(X_j,Y_j)$ have unit variances
and correlation $0<r_0<1$, and put
$O=a^{-2}\sum_{j=1}^M(X_j^2-1)$. Then
\begin{equation}
\E(O\mid Y)=a^{-2}r_0^2\sum_j(Y_j^2-1),\quad
\E\Var(O\mid Y)=2a^{-4}M(1-r_0^4),\quad
\Var(\E(O\mid Y))=2a^{-4}Mr_0^4.
\label{eq:Gaussian-counterexample}
\end{equation}
For $M\asymp a^{-4}$ both pieces have order $a^{-8}$. Thus the exact subtraction
isolates a genuine ultraviolet contact but does not itself prove tightness of the
remaining observable. This test is a separate probability model, not a replacement
of the nonlinear Wilson theory.

Finite translation/rotation averages of the new hierarchies are convex mixtures
of projected-observable laws, not the law of the averaged observable. They retain
degree-two contact locality. At higher order, the logarithm of a global mixture
can correlate distinct patches; an all-order local symmetric prescription does
not follow from that averaging alone. Reflection positivity follows from the
original law only when each chosen scheme commutes with the same reflection and
the projected positive-time observables remain in its positive half. Pairing
two reflected schemes is not by itself that proof.

\section{A separate Gaussian contact extension and a positivity obstruction}
\label{sec:Gaussian}

This section is an explicit test in the free linearized theory. It does not infer
interacting estimates from perturbation theory. On a fixed physical torus define
the Gaussian law in Feynman gauge: its color/component Fourier covariance is
$\delta_{ab}\delta_{\mu\nu}/\widehat p^2$ for $p\ne0$, up to a specified overall
positive normalization, with
$\widehat p^2=\sum_\mu4a^{-2}\sin^2(ap_\mu/2)$; the zero mode is removed.
Equivalently one may use the transverse inverse, since the currents annihilate
longitudinal modes. Retain the original forward/backward difference phases.
The linear directional currents
$j_{\nu;\mu}=d_\mu^*F_{\mu\nu}$ have two derivatives of the gauge potential.
For their covariance matrix $C_a(z)$,
\begin{equation}
|C_a(z)|\le C_T(d_T(z,0)+a)^{-6},
\qquad C_a\longrightarrow C\quad\hbox{smoothly off }0.
\label{eq:freecurrent}
\end{equation}
A direct Fourier proof decomposes the multiplier $D^4/\widehat p^2$ into
smooth dyadic shells $|p|\asymp K$. Its $r$th derivatives are bounded by
$C_rK^{2-r}$, so summation by parts gives
$C_NK^6(1+K|z|)^{-N}$ for a shell kernel. Summing to $K\asymp a^{-1}$ gives
\eqref{eq:freecurrent}; periodization gives the fixed-torus bound. The zero mode
is killed by the differences. Fixed pointed shifts are $O(a)$ and change only
the constants in the estimate.

Let $P$ be the block projection $I_3-\frac13\mathbf1\mathbf1^T$ on plane slots
for each $\nu$. For an unshifted representative the centered quadratic density
has covariance
\begin{equation}
C_{H,a}(z)=2m\Tr\{P C_a(z)P C_a(-z)\},
\qquad |C_{H,a}(z)|\le C_T(d_T(z,0)+a)^{-12}.
\label{eq:freeH}
\end{equation}
The completed observable uses its literal finite cross-average, which has the
same bound and off-diagonal limit. This is Wick's identity in the separate
Gaussian model. A single-source identity subtraction, of order $a^{-6}$, is
distinct from the two-source prescription below.

\begin{theorem}[Explicit two-point extension]
\label{thm:extension}
Choose a smooth cutoff $\chi$ supported in a coordinate ball and equal to one
near zero. For a smooth test $\varphi$ put
\begin{align}
W_8\varphi(z)&=\varphi(z)-\chi(z)
\sum_{|\alpha|\le8}\frac{z^\alpha}{\alpha!}\partial^\alpha\varphi(0),\label{eq:W8}\\
\mathcal R_a(\varphi)&=a^4\sum_z C_{H,a}(z)W_8\varphi(z).
\label{eq:Ra}
\end{align}
Then $\mathcal R_a$ converges as a distribution to
\begin{equation}
\mathcal R(\varphi)=\int_{z\ne0}C_H(z)W_8\varphi(z)\dd z,
\label{eq:Rlimit}
\end{equation}
with a uniform $C^9$ test bound on fixed compact sets. The finite-cutoff
counterterm coefficients are explicitly
\begin{equation}
c_\alpha(a)=a^4\sum_z C_{H,a}(z)\chi(z)\frac{z^\alpha}{\alpha!},
\quad |c_\alpha(a)|\le
\begin{cases}
C a^{|\alpha|-8},&|\alpha|<8,\\
C[1+\log(1/a)],&|\alpha|=8.
\end{cases}
\label{eq:countercoeff}
\end{equation}
Fixed finite local terms can be added to impose specified finite normalizations.
\end{theorem}
\begin{proof}
Near zero, Taylor's theorem gives
$|W_8\varphi(z)|\le C\|\varphi\|_{C^9}|z|^9$.
The absolute contribution from $|z|<\delta$ is bounded uniformly in $a$ by
$C\|\varphi\|_{C^9}\delta$: in four dimensions the radial power is
$r^3r^{-12}r^9=r^0$. The $z=0$ term vanishes exactly. Away from the ball,
\eqref{eq:freecurrent} gives ordinary Riemann-sum convergence. Sending first
$a\to0$ and then $\delta\to0$ proves the claim. The same radial counting with
$|z|^{|\alpha|}$ proves \eqref{eq:countercoeff}.
For smeared two-point tests use
$\varphi(z)=\int f(x)h(x+z)\dd x$ and its lattice version; their $C^9$
convergence is sufficient by the proved uniform bound.
\end{proof}

The general scaling-degree extension framework is described in
Brunetti--Fredenhagen, \S5~\cite{bf}; the elementary proof above establishes this
particular two-point assertion directly. It is not an interacting construction.

For a hypercubic scalar and a hypercubic cutoff, coordinate-odd coefficients
vanish. With $p_{2k}=\sum_\mu\partial_\mu^{2k}$, the local ambiguity through
order eight has the following twelve independent invariant structures:
\begin{center}
\begin{tabular}{cl}
\toprule Differential order&Invariant operators acting on $\delta$\\\midrule
0&$1$\\
2&$p_2$\\
4&$p_2^2,\ p_4$\\
6&$p_2^3,\ p_2p_4,\ p_6$\\
8&$p_2^4,\ p_2^2p_4,\ p_4^2,\ p_2p_6,\ p_8$\\\bottomrule
\end{tabular}
\end{center}
These are symmetric polynomials in the four squared derivative variables through
degree four. Full rotational symmetry would impose additional restrictions;
it is not assumed for this hypercubic $H$ operator.

\subsection{The logarithmic residue can be computed explicitly}

On $\R^4$ fix the free normalization
\begin{equation}
\E A_\mu^a(x)A_\nu^b(0)=\delta_{ab}\delta_{\mu\nu}c|x|^{-2},\qquad
j_{\nu;\mu}=-\partial_\mu^2A_\nu+\partial_\mu\partial_\nu A_\mu,
\qquad H=\sum_\nu j_\nu^TPj_\nu.
\label{eq:freeconvention}
\end{equation}
Away from zero its covariance is $C_H(x)=|x|^{-12}\mathcal A(x/|x|)$,
with $\mathcal A\ge0$ and not identically zero. Positivity follows from
$2m\|P^{1/2}C(x)P^{1/2}\|_{\rm HS}^2$. The covariance of the total current
with a directional current has polynomial Fourier multiplier and is supported
at zero; hence $PC(x)P=C(x)$ away from zero. For example, at $x=(r,0,0,0)$,
the $j_{2;1}$ self-covariance is $80c r^{-6}$, proving nonvanishing.

Writing $q_{2k}(\omega)=\sum_\mu\omega_\mu^{2k}$ on $S^3$, direct differentiation
and Wick contraction give
\begin{equation}
\frac{\mathcal A(\omega)}{2mc^2}
=-4608+124416q_4+294912q_4^2-368640q_6.
\label{eq:angular}
\end{equation}
The formula follows by using
\begin{align*}
\partial_i\partial_j\partial_k\partial_l|x|^{-2}
={}&8|x|^{-6}(\delta_{ij}\delta_{kl}+\delta_{ik}\delta_{jl}+\delta_{il}\delta_{jk})\\
&-48|x|^{-8}\sum_{\text{six pairs}}\delta_{ij}x_kx_l
+384|x|^{-10}x_ix_jx_kx_l
\end{align*}
in the full $12\times12$ directional-current covariance matrix and summing its squared entries.
An exact-arithmetic check verifies the polynomial identity, every
total-current cancellation, and every degree-eight coefficient below.

Define now $p_{2k}=\sum_\mu p_\mu^{2k}$ for a Fourier variable $p$, and let
\begin{equation}
\mathcal P_8(p)=\frac1{8!}\int_{S^3}\mathcal A(\omega)(p\cdot\omega)^8\dd\omega.
\label{eq:P8integral}
\end{equation}
The exact sphere moments give
\begin{equation}
\boxed{\mathcal P_8(p)=\frac{\pi^2mc^2}{2520}
\left(297p_2^4+96p_2^2p_4+32p_4^2-104p_2p_6\right).}
\label{eq:P8explicit}
\end{equation}
For $c=(4\pi^2)^{-1}$ the prefactor is $m/(40320\pi^2)$.
Since $p_6\le p_2^3$, the bracket is at least $193p_2^4$; in particular the
residue is strictly positive for $p\ne0$. The sphere-moment calculation uses
\[
\frac1{|S^3|}\int_{S^3}\omega^\alpha\dd\omega
=\frac{\prod_\mu(\alpha_\mu-1)!!}
{4\cdot6\cdots(4+|\alpha|-2)}
\]
when every exponent is even, and zero otherwise. Together with
\eqref{eq:angular} this is a finite coefficientwise derivation of
\eqref{eq:P8explicit}, rather than numerical quadrature.

\begin{theorem}[Ordinary covariance positivity fails at the allowed contact order]
\label{thm:positivity-obstruction}
Let $T$ be any translation-invariant extension of the free $\R^4$ kernel
$C_H(x)$ with scaling degree at most twelve. No such extension is a
positive-definite ordinary Euclidean covariance. More explicitly, with Fourier
convention $\widehat T(p)=\langle T,e^{-ip\cdot x}\rangle$,
\begin{equation}
\Re\widehat T(\lambda p)
=-\lambda^8\log\lambda\,\mathcal P_8(p)+O_p(\lambda^8),
\qquad\lambda\to\infty,\quad p\ne0.
\label{eq:negativeFourier}
\end{equation}
\end{theorem}
\begin{proof}
Every allowed extension differs from the radial-cutoff Taylor extension by
derivatives of $\delta$ of order at most eight, hence by a Fourier polynomial of
degree at most eight. For the Taylor extension its Fourier transform is a continuous
function given by the convergent subtracted oscillatory integral. Dilate
$x\mapsto x/\lambda$. The degree-eight Taylor term contributes
\[
\lambda^8\mathcal P_8(p)
\int_0^\infty[\chi(r)-\chi(r/\lambda)]\frac{\dd r}{r}
=-\lambda^8\log\lambda\,\mathcal P_8(p),
\]
because $(-i)^8=1$. Every lower Taylor term and every finite allowed contact
change is $O_p(\lambda^8)$. This proves \eqref{eq:negativeFourier}.
Its real part is negative on a neighborhood of every sufficiently large point
along any fixed nonzero ray. A Schwartz test whose Fourier support lies there
has negative ordinary covariance pairing. Thus positive definiteness fails.
For a real even extension its Fourier transform is real; taking the real part
also covers extensions with odd contact terms.
\end{proof}

This obstruction concerns ordinary probability covariance. It does not contradict
Osterwalder--Schrader reflection positivity on separated time supports, where the
local counterterms vanish from the reflected pairing. It is proved for the free
$\R^4$ model and for extensions respecting the stated scaling degree; higher-order
contact terms are outside the theorem. A full hierarchy still requires compatible
extensions at every partial diagonal. Fixed positive flow-time smoothing does not
by itself settle that local hierarchy.

The restriction on contact order matters. The Fourier transform of the real even
Taylor extension is bounded in absolute value by
$C(1+|p|)^8\log(2+|p|)$. Adding a sufficiently large multiple of
$\delta+(-\Delta)^5\delta$ therefore makes its Fourier transform nonnegative.
This changes the scaling degree to fourteen. It repairs this two-point ordinary
positivity test at the cost of a higher-order contact. The resulting positive
covariance can define a Gaussian hierarchy, but this does not construct the
specified composite-$H$ higher-point hierarchy or the interacting hierarchy.

There is also a useful diagnostic for the prescription of
Section~\ref{sec:conditional}. In the separate free model, a sequence of
ordinary positive covariances cannot have a translation-invariant
distributional limit on $\R^4$ with the same off-diagonal kernel
$C_H$ and scaling degree at most twelve: distributional convergence preserves
positive definiteness, contradicting Theorem~\ref{thm:positivity-obstruction}.
Consequently a proposed positive projected-observable limit with that
off-diagonal kernel must either fail to converge or leave this contact-order
class. This is not a no-go theorem for the interacting source, whose necessary
off-diagonal comparison has not been established.

\subsection{The exact interacting input still needed}

The same Taylor argument would apply to an actual interacting connected kernel
after specified single-operator mixing if one proves off-diagonal convergence and
\begin{equation}
|C_{H,a}^{\rm int}(z)|\le
C(|z|+a)^{-12}\bigl(1+|\log(|z|+a)|\bigr)^p
\label{eq:interactingneeded}
\end{equation}
for a fixed $p$, uniformly on the required volumes. The near-diagonal remainder
would be $O(\delta(1+|\log\delta|)^p)$, and the same contact order would suffice.
Neither this interacting bound nor its all-order partial-diagonal analogues has
been proved here. The exact nonlinear prescription in Section~\ref{sec:conditional}
and the separate Gaussian extension in this section are therefore kept distinct.

\section{Actual one-step selector admission and current-coordinate lifts}
\label{sec:mixed}

The map in this section is the original one-step average $q:\mathcal D\subset
G^E\to G^C$ from ME14--31. It is not the matched-reference map of a general
$R$ history. Its complete row readset $S(c)$ and central crossing pivot $p(c)$
have the exact properties
\begin{equation}
S(c)\cap\{p(d):d\in C\}=\{p(c)\},\qquad
dq_c(X_{p(c),\cdot})\big|_{U=I}=L^{-3}I.
\label{eq:pivotinput}
\end{equation}
The first property includes every original rooted contour and central path.
For the second, differentiating the original reference-logarithm mean cancels
the two reference-path terms; exactly $L$ of its $L^4$ translated middle paths
read the crossing pivot. Endpoint means stay in their endpoint blocks. The
$L^{-3}$ is a group/Haar-coordinate factor; physical connection coordinates
have an additional length conversion.

Write $dq_c(V)=(Vq_c)q_c^{-1}$ and
\begin{equation}
A_c=dq_c(X_{p(c),\cdot}),\qquad B_{ci}=dq_c(X_i).
\label{eq:AB}
\end{equation}

\begin{theorem}[Literal local plaquette admission]
\label{thm:admission}
Suppose that for every row touched by a source direction its full original readset
lies in a contractible current-lattice box $Q_c$ of side at most $K=O(L)$ bonds,
with no omitted winding or prescribed seam. Suppose the actual selector implies
$\|U(\partial p)-I\|\le\epsilon$ on every plaquette of that box. There is a fixed
$\epsilon_L>0$ such that, if $3K\epsilon\le\epsilon_L$, the actual row belongs
to its pivot chart and
\begin{equation}
s_{\min}(A_c)\ge\tfrac12L^{-3},\qquad
s_{\max}(A_c)\le\tfrac32L^{-3},\qquad
\|A_c^{-1}\|\le2L^3.
\label{eq:pivotbound}
\end{equation}
Every fixed-order original-coordinate jet of $A_c^{-1}$ and $B_{ci}$ has a
bound $C_{L,r}$ on a fixed smaller chart. These constants do not grow with the
number of active rows or with an earlier ultraviolet mesh.
\end{theorem}
\begin{proof}
Use the coordinate-ordered tree rooted at a box corner. A tree fundamental loop
is filled by commuting its final unit step past the later coordinate segments;
at most $3K$ original plaquettes are needed. Its holonomy is the corresponding
product of transported plaquettes. The transports are unitary on the real field,
so telescoping gives $\|U_b^g-I\|\le3K\epsilon$ for every bond in the box.
All row and endpoint frames transform together. Analyticity of the original
finite mean and \eqref{eq:pivotinput} now give
$\|L^3A_c-I\|\le1/2$ in transported frames. Fixed real adjoint transports are
isometries, giving the intrinsic singular-value bounds \eqref{eq:pivotbound}.

For derivative estimates one can hold the real center gauge fixed. The original
paths have bounded length at fixed $L$; in a fixed smaller complex tube their
products and inverses remain uniformly bounded and the mean remains on its
actual branch. Finite-variable Cauchy estimates and inversion give $C_{L,r}$.
This argument uses a local gauge to prove a coordinate estimate, without changing
the original integration variable or erasing exterior data.
\end{proof}

For the actual source bond support $S$, define
\begin{equation}
C_S=\{c:S(c)\cap S\ne\varnothing\},\qquad P_S=\{p(c):c\in C_S\}.
\label{eq:activerows}
\end{equation}
The one-step geometry gives bounded row and column overlap:
$|C_S|+|\bigcup_{c\in C_S}S(c)|\le C_L|S|$. By pivot exclusion the active
pivot derivative is block diagonal. Thus the inverse function theorem supplies
the actual simultaneous local chart
\begin{equation}
\Phi_S(U)=(z=q_{C_S}(U),w=U_{E\setminus P_S}),\qquad U=\Psi_S(z,w).
\label{eq:partialchart}
\end{equation}
Every other original row remains its actual function of $w$. Rough unrelated
exterior fields are not put in a potential chart; the original averages that
are used must still lie on their declared branches.

\begin{proposition}[Complete tangent and output decomposition]
\label{prop:lifts}
For a nonpivot source direction define
\begin{align}
a_{ci}&=A_c^{-1}B_{ci},
&H_i&=X_i-\sum_{c,\alpha}a_{ci}^\alpha X_{p(c),\alpha},\notag\\
Z_{c,\alpha}&=\sum_\beta(A_c^{-1})^{\beta\alpha}X_{p(c),\beta}.
\label{eq:lifts}
\end{align}
Then
\begin{align}
dq_d(H_i)&=0\quad\hbox{for every original one-step row }d,\notag\\
(\Phi_S)_*H_i&=X_{w_i},
&(\Phi_S)_*Z_{c,\alpha}&=X_{z_c,\alpha},\notag\\
X_i&=\ind_{\{i\text{ nonpivot}\}}H_i
+\sum_{c,\alpha}B_{ci}^\alpha Z_{c,\alpha}.
\label{eq:complete-lift}
\end{align}
For $V=\sum_{i\in S}v_iX_i$, this gives the exact split
\begin{equation}
V=V_{\rm fib}+V_{\rm out},\quad
V_{\rm fib}=\sum_{i\text{ nonpivot}}v_iH_i,\quad
V_{\rm out}=\sum_{c,\alpha}(dq_cV)^\alpha Z_{c,\alpha}.
\label{eq:Vsplit}
\end{equation}
\end{proposition}
\begin{proof}
For an active row the compensation is exactly $B_{di}-A_dA_d^{-1}B_{di}=0$.
A remote row reads neither the original source direction nor an active pivot,
by \eqref{eq:pivotinput}; no inverse of a rough remote row is needed. The $w$
coordinates change exactly by their original invariant field. The output
identities and the pivot case follow directly from \eqref{eq:AB}.
\end{proof}

Every source pivot has its full output response and zero internal component in
this chart. The coefficient projection in \eqref{eq:Vsplit} has bounded
$\ell^2\to\ell^2$ norm by its two bounded Schur sums. Fixed ordered words are
transformed in their original order, including all derivatives of $A_c^{-1}$ and
$B_{ci}$. For the convention of Section~\ref{sec:conventions},
$[H_{e,A},H_{f,B}]=-\delta_{ef}f_{AB}^{\ C}H_{e,C}$ when
$[T_A,T_B]=f_{AB}^{\ C}T_C$; the coordinate pushforward preserves this sign.
No all-history contraction follows from the finite one-step jet bound.

The vectors in \eqref{eq:lifts} are intrinsic functions of the same original map
and agree on overlaps of its admitted charts. If a genuine chart partition
$\sum_\alpha\psi_\alpha=1$ is used near a distribution's support, linearity gives
\begin{equation}
\sum_\alpha D_{H_{i_r}}\cdots D_{H_{i_1}}(\psi_\alpha T)
=D_{H_{i_r}}\cdots D_{H_{i_1}}T.
\label{eq:chartglue}
\end{equation}
Thus artificial chart-interface derivatives cancel before estimates. This does
not cancel original history selectors or compare different history-dependent maps.

\subsection{Coarea, full effective scores, and sharp selectors}

The exact Haar factor is
\begin{equation}
\Psi_S^*\dd\lambda_U=J_S\dd\lambda_z\dd\lambda_w,
\qquad J_S(z,w)=\prod_{c\in C_S}
|\det A_c(\Psi_S(z,w))|^{-1}.
\label{eq:coarea}
\end{equation}
The crucial divergence sign is
\begin{equation}
X_{w_i}\log J_S=(\diver H_i)\circ\Psi_S,
\qquad\diver H_i=-\sum_{c,\alpha}X_{p(c),\alpha}a_{ci}^\alpha.
\label{eq:coareascore}
\end{equation}
Indeed the coordinate invariant field has zero divergence for product Haar,
and divergence for $J_S\lambda_z\lambda_w$ adds $X_{w_i}\log J_S$.
Logarithmic coordinates must also retain their Haar coordinate densities.

For the complete positive smooth body $\rho=e^{-\beta_sS_0}G$,
$R=(\rho\circ\Psi_S)J_S$, the negative internal score is exactly
\begin{equation}
-X_{w_i}\log R
=\bigl[\beta_sH_iS_0-H_i\log G-\diver H_i\bigr]\circ\Psi_S.
\label{eq:fullscore}
\end{equation}
A variable-coefficient vector uses
$\rho^{-1}D_{V_{\rm fib}}\rho
=\beta_sV_{\rm fib}S_0-V_{\rm fib}\log G-\diver V_{\rm fib}$;
the divergence includes $\sum_iH_iv_i$. For signed or vanishing selected terms
use the unnormalized tested identity, without inventing a logarithm or a positive
conditional law:
\begin{equation}
(\Phi_S)_*(D_{X_i}\rho)
=-\ind_{\{i\text{ nonpivot}\}}X_{w_i}R
-\sum_{c,\alpha}X_{z_c,\alpha}\bigl[(B_{ci}^\alpha\circ\Psi_S)R\bigr].
\label{eq:pushcurrent}
\end{equation}
This retains derivatives of the coarea factor and the moving output coefficients.

For the smooth Wilson body alone, compact plaquette estimates give, with literal
local star unions $Y_i$ and $E_Y=\beta_s\sum_{p\in Y}s_p$,
\begin{align}
|-X_{w_i}\log R|&\le C_L(1+\sqrt{\beta_s}\sqrt{E_{Y_i}}),\notag\\
|X_{w_{i_r}}\cdots X_{w_{i_1}}R|
&\le C_{L,r}\beta_s^{r/2}(1+E_{Y_r})^{r/2}R,\qquad\beta_s\ge1.
\label{eq:scorejets}
\end{align}
Higher log derivatives are bounded by $C_{L,r}(1+\beta_s)$; ordered exponential
differentiation proves the second estimate. A full inherited $G$ adds its actual
ordered jets $H_{i_j}\cdots H_{i_1}\log G$ to the finite Bell polynomial.
Normalization of $G$ supplies no estimate for these derivatives.

Let $\Xi$ denote the original joint sharp selectors, gauge factors, and old
constraints. It remains $\Xi\circ\Psi_S$. For a smooth multiplier $\chi$,
$D_i=-X_{w_i}$ gives
\begin{equation}
[D_r\cdots D_1,\chi]
=\sum_{j=1}^r D_r\cdots D_{j+1}M_{D_j\chi}D_{j-1}\cdots D_1.
\label{eq:boundarycommutator}
\end{equation}
For a sharp selector this is distributional; already singular constraints are
differentiated as their complete original joint distribution, without forming
undefined products of exposed delta factors. If an actual selector is
$\chi(U;b)=\chi_{\rm out}(q(U);b)$ with retained reference data satisfying
$H_i b=0$, then every checked $H_i\chi=0$ and all
ordered internal words commute with it. Its boundary response remains in
$V_{\rm out}$. Fine selectors and older differently averaged constraints need
their own check.

The effect of replacing the bare score by the full score is especially explicit:
\begin{equation}
O_H\Xi G\rho_W
=r_H\Xi G\rho_W+D_{V_H}(\Xi G\rho_W)+\rho_WV_H(\Xi G).
\label{eq:historycorrection}
\end{equation}
The last term contains $+V_H\log G$ in its smooth ordinary part and every original
selector/constraint response. Keeping only the old $r_H$ after changing the score
would violate this identity.

\subsection{Actual $H$ projection, retained constraints, and live references}

On the lattice where the original source is evaluated, keep
$\sigma=\gamma^{-2}a^{-2}$ and its action coefficient $\beta_s$. The full
physical vector has the internal component
\begin{equation}
V_{H,\rm fib}(f)=-\frac{\sigma}{3\beta_s}
\sum_{i\text{ nonpivot}}f_iF_iH_i.
\label{eq:Hfib}
\end{equation}
Let $Y$ be the literal union of source and compensating-pivot plaquette stars,
and $N_S$ the actual number of source bond/color directions. Bounded overlap and
$|F_i|^2\le C\sum_{p\ni i}s_p$ prove
\begin{align}
\|V_{H,\rm fib}\|_{\rm coeff}^2
&\le C_L\sigma^2\beta_s^{-3}\|f\|_\infty^2 E_Y,\label{eq:Hfibbound}\\
|t_{H,\rm fib}|
&\le C_L|\sigma|\beta_s^{-1}\|f\|_\infty(N_S+E_Y)
\quad(G=1,\ \beta_s\ge1).
\label{eq:Hscorebound}
\end{align}
To verify the second bound, the action term is a finite-overlap product of two
forces and costs $C_L|\sigma|\beta_s^{-1}\|f\|_\infty E_Y$.
Differentiating $f_iF_i$ contributes bounded second Wilson derivatives times
$N_S$. The remaining coefficient-divergence term is bounded by the same prefactor
times $\sqrt{N_SE_Y/\beta_s}$ and is absorbed by $N_S+E_Y$.
The complete inherited body adds exactly $-V_{H,\rm fib}\log G$ and the original
sharp distributional partners. If $\beta_s=c_\beta\gamma^{-2}$, then
$\sigma\beta_s^{-1}=c_\beta^{-1}a^{-2}$: no missing birth gain is supplied.
Conversion to $U=e^{hB}U^0$ requires the actual inverse exponential differential
and a factor $h^{-1}$, hence $h^{-2}$ in the squared vector norm. An identification
of this fiber with the native $B$ fiber is not assumed.

The retained-halo chart ME19--31 has
$(h,z,w)=(U_{H^+},q_{O_H},w)$, with
$q=(q_{I_H}(h),z)$. At a prescribed full output it still contains
\begin{equation}
\delta(q_{I_H}(h)-y_{I_H}).
\label{eq:remainingdelta}
\end{equation}
A retained inside pivot has the actual flat derivative
$X_{p(c),A}q_c=L^{-3}T_A$; differentiating \eqref{eq:remainingdelta} therefore
produces a nonzero output-delta response. Solving the inside pivots recovers the
full-output tangent chart and the exact inside determinant factors. A retained
direction at fixed $z$ is not automatically internal at fixed full $q$.

For a parameter family of this same original one-step map, preserving its fixed
read sets and pivot exclusion, let $b$ be a live reference and $y_c(b)$ a moving
group target. Use the relative output $r_c=q_c(U;b)y_c(b)^{-1}$. Define
$u_c=(\partial_bq_c)q_c^{-1}$ and $v_c=(\partial_by_c)y_c^{-1}$.
The exact fixed-relative-output lift is
\begin{equation}
R_b^{\rm rel}=\partial_b-\sum_{c,\alpha}
\left[A_c^{-1}(u_c-\Ad_{r_c}v_c)\right]^\alpha X_{p(c),\alpha}.
\label{eq:movingtarget}
\end{equation}
The sum includes every row moved by either $q$ or $y$, including any outside
$C_S$, with admission verified on each such row. Only on $q_c=y_c$ can the
adjoint be removed. The rowwise formula does not apply to a general mixed map
with off-diagonal pivot derivatives. Old prefixes and reference derivatives
acquire no automatic $C_L$ bound. The reference base measure, coarea, all
independent scalar derivatives, and the responses of every other
$b$-dependent selector, gauge factor, or domain indicator remain in the full
density divergence.

\begin{remark}
This section closes an actual one-step current-coordinate identification on a
specified local selector sector. It does not identify arbitrary $R$-I/$R$-II
prepared fields or old readouts with that sector, prove the three localization
joins of AE8/LM47, or give a spectral gap after additional selectors. An admitted
coordinate chart can still contain a disconnected selected support.
\end{remark}

\section{Exact complete-score transport on a verified positive fiber}
\label{sec:Poisson}

Fix all retained old variables and live references $b$. On an \emph{identified}
native fiber let
\begin{equation}
\dd\mu_b=\rho_b(B)\ind_{\Omega_R}(B)\dd B,\qquad
Z_b=\mu_b(1)>0,\qquad K_b=Z_b^{-1}\mu_b.
\label{eq:actualfiber}
\end{equation}
The complete smooth positive body $\rho_b$ contains every inherited effective
factor and Jacobian. The sharp product-norm cutoff is the original one. Assume
its actual Poincar\'e inequality
\begin{equation}
\Var_{K_b}F\le C_P(b)K_b|\nabla_BF|^2.
\label{eq:actualgap}
\end{equation}
Reference~\cite[\S4]{ND} verifies this for its admitted reference and small same-record weight
interpolation. If $\rho_b=f_be^{-W_b}$ on that same support and the reference
constant is $C_0$, a direct density-ratio comparison gives
\begin{equation}
C_P(b)\le e^{\osc_{\Omega_R}W_b}C_0.
\label{eq:oscillation}
\end{equation}
Indeed $r=\dd K_b/\dd K_0$ gives
$\Var_{K_b}F\le(\sup r)\Var_{K_0}F
\le C_0(\sup r/\inf r)K_b|\nabla F|^2$.
Large constant parts of $W_b$ remain in $Z_b$. A new sharp selector is not a
bounded-oscillation perturbation: a Gaussian restricted to two disjoint intervals
has a nonconstant piecewise-constant function with zero internal gradient.
No spectral gap is imported into the one-step chart of Section~\ref{sec:mixed}.

Define the full effective-score operator
\begin{equation}
L_b=\Delta_B+\nabla\log\rho_b\cdot\nabla_B.
\label{eq:Lb}
\end{equation}
For $KF=0$ and $F\in L^2(K)$, the mean-zero weak Neumann problem is
\begin{equation}
K(\nabla u_F\cdot\nabla\phi)=K(F\phi)\quad(\phi\in H^1(K)),
\qquad Ku_F=0.
\label{eq:weakPoisson}
\end{equation}
Coercivity \eqref{eq:actualgap} and the Hilbert-space representation theorem give
a unique solution. Set $V_F=\nabla u_F$. In weak density notation,
\begin{equation}
D_{V_F}\mu=F\mu,\qquad
\|V_F\|_2\le\sqrt{C_P}\|F\|_2,\qquad
\|V_F\|_2\le C_P\|\nabla F\|_2\quad(F\in H^1).
\label{eq:inverse}
\end{equation}
The first bound follows by testing with $u_F$ and applying Poincar\'e; the second
then applies Poincar\'e to $F$. The reflecting condition is imposed on this new
auxiliary current. It makes no assertion about the flux of an old physical current.

\begin{theorem}[Finite-source transport with the original mass retained]
\label{thm:Poissontransport}
For the actual source action $A(\eta,B)$ define
\begin{equation}
c(\eta;b)=\log K_b e^A,\qquad R(\eta,B;b)=e^{A-c}-1.
\label{eq:centeredsource}
\end{equation}
Every coefficient $R_I$ has $K_bR_I=0$. If these coefficients lie in $L^2(K_b)$,
let $V_I=\nabla(-L_b)^{-1}R_I$ in the sense of \eqref{eq:weakPoisson}. Then
\begin{align}
\mu_be^A&=e^c\left[\mu_b+\sum_{I\ne\varnothing}\eta_I D_{V_I}\mu_b\right],
\label{eq:transportdensity}\\
\int e^A\Phi\dd\mu_b
&=Z_be^c\left[K_b\Phi+\sum_{I\ne\varnothing}\eta_I
K_b(V_I\cdot\nabla_B\Phi)\right].
\label{eq:transportcarrier}
\end{align}
The tested current formula holds coefficientwise for $H^1$ carrying functions.
\end{theorem}
\begin{proof}
$Ke^{A-c}=1$ in the finite algebra, so each $R_I$ is centered. Substitute
\eqref{eq:inverse} into $\mu_be^{A-c}=\mu_b(1+R)$ and multiply by the
internally constant $e^c$. This requires no composition or exponential flow of
weak vector fields. Testing gives \eqref{eq:transportcarrier}.
\end{proof}

Every generated scalar/contact remains in $e^c$. Deleting that factor would
change the original physical source. If the original expression is
$e^{s(b)}D_{\rm orig}(b)^{-1}\int\rho_bF$, its scalar is
\begin{equation}
\frac{e^{s(b)}Z_b}{D_{\rm orig}(b)}K_bF;
\label{eq:absolutescalar}
\end{equation}
the ratio is not set to one, and every retained derivative differentiates it.
For an actual readout $q(b,B)$,
$V_I\cdot\nabla(\phi\circ q)=(D_Bq\,V_I)\cdot\nabla\phi$.
The absolute mass of the pushed vector is at most
$Z_b\|\|D_Bq\|_{\rm op}\|_{L^2(K_b)}\|V_I\|_2$, with all actual coordinate components.
These currents can depend on the whole native component; no spatial-distance
gain or retained-variable derivative bound follows from \eqref{eq:inverse}.

\subsection{Two-source currents and a quantitative $L^4$ bound}

For \eqref{eq:source}, let $m_i=KO_i$ and $t_i^P=O_i-m_i$. Exact finite algebra gives
\begin{equation}
c_{12}=KA_{12}+\Cov_K(O_1,O_2),\qquad
R_{12}=A_{12}+t_1^Pt_2^P-c_{12}.
\label{eq:actualc12}
\end{equation}
An ordinary representative of its covariance part is
\begin{equation}
\Gamma_{12}^P=\tfrac12(V_{t_1^P}\cdot\nabla O_2+V_{t_2^P}\cdot\nabla O_1),
\quad K\Gamma_{12}^P=\Cov_K(O_1,O_2),
\quad K|\Gamma_{12}^P|\le C_P\|\nabla O_1\|_2\|\nabla O_2\|_2.
\label{eq:contactrepresentative}
\end{equation}
It represents the actual nonzero contact. Its raw mass includes $Z_b$ and the
absolute value of every independent prefactor.

Poincar\'e applied to $(F-KF)^2$ supplies a useful fourth-moment consequence:
\begin{equation}
\|F-KF\|_4\le c_4\sqrt{C_P}\|\nabla F\|_4,
\qquad c_4=\sqrt{2+\sqrt5}.
\label{eq:L4Poincare}
\end{equation}
In fact, with $X=(K(F-KF)^4)^{1/2}$ and $a_4=\|\nabla F\|_4^2$, the
inequality is $X^2\le C_P^2a_4^2+4C_Pa_4X$. Solving the quadratic proves
\eqref{eq:L4Poincare}; truncation gives the usual Sobolev version.
Since $\nabla R_{12}=\nabla A_{12}+t_1^P\nabla O_2+t_2^P\nabla O_1$,
\eqref{eq:inverse}, H\"older, and \eqref{eq:L4Poincare} give
\begin{equation}
\boxed{\|V_{12}\|_2\le C_P\|\nabla A_{12}\|_2
+2c_4C_P^{3/2}\|\nabla O_1\|_4\|\nabla O_2\|_4.}
\label{eq:V12bound}
\end{equation}
Large insertion means are absent from this bound. The full effective score is
still present in the inverse defining $V_{12}$.
At degree two an external carrier also contributes
$K(V_1\cdot\nabla\Phi_2)+K(V_2\cdot\nabla\Phi_1)
+K(V_{12}\cdot\nabla\Phi_0)$, followed by the full convolution with $e^c$.
The scalar coefficient $c_{12}$ alone does not describe that carrying response.

\subsection{Exact conversion of the old boundary and Casimir partners}

Suppose an old decomposition has been correctly lifted to this same realization:
\begin{equation}
O_i=t_i+r_i,\qquad t_i\mu=D_{V_i^{\rm old}}\mu+B_i.
\label{eq:oldboundarysplit}
\end{equation}
$B_i$ is the complete original boundary/constraint partner, with its original
sign. For a smooth old field on the fixed cutoff,
\begin{equation}
t_i=-\rho^{-1}\diver(\rho V_i^{\rm old}),\qquad
B_i=-\rho(V_i^{\rm old}\cdot n)\dd\sigma_{\partial\Omega_R}.
\label{eq:boundarysign}
\end{equation}
The sign follows from $\nabla\ind_{\Omega_R}=-n\delta_{\partial\Omega_R}$.
The new Poisson current satisfies the exact difference identity
\begin{equation}
\boxed{D_{V_i^P-V_i^{\rm old}}\mu=(r_i-m_i)\mu+B_i.}
\label{eq:changeofsplit}
\end{equation}
Thus its zero flux does not remove the old partner.

Let $\mathcal B_i(F)=Z^{-1}\langle B_i,F\rangle$ and
$\mathring{\mathcal B}_i(F)=\mathcal B_i(F)-\mathcal B_i(1)KF$.
Retain $b_{12}^{\rm old}$ from \eqref{eq:oldb}, with these actual $t_i,r_i$.
Then
\begin{equation}
\boxed{\begin{aligned}
c_{12}={}&Kb_{12}^{\rm old}+\Cov_K(r_1,r_2)\\
&+\mathring{\mathcal B}_1(r_2+\tfrac12t_2)
+\mathring{\mathcal B}_2(r_1+\tfrac12t_1).
\end{aligned}}
\label{eq:boundarybridge}
\end{equation}
To prove it, use
$K(t_iF)=K(V_i^{\rm old}F)+\mathcal B_i(F)$, expand $O_1O_2$, symmetrize
the pure-current term, and subtract the full product of means. The centered
boundary terms include $-\mathcal B_1(1)\mathcal B_2(1)$.

For external carriers, the raw ordered identity is more detailed:
\begin{align}
(A_{12}+O_1O_2)\mu={}&(b_{12}^{\rm old}+r_1r_2)\mu
+D_{V_1^{\rm old}}(r_2\mu)+D_{V_2^{\rm old}}(r_1\mu)\notag\\
&+\tfrac12(D_{V_1^{\rm old}}D_{V_2^{\rm old}}
+D_{V_2^{\rm old}}D_{V_1^{\rm old}})\mu\notag\\
&+(r_2+\tfrac12t_2)B_1+(r_1+\tfrac12t_1)B_2\notag\\
&+\tfrac12(D_{V_1^{\rm old}}B_2+D_{V_2^{\rm old}}B_1).
\label{eq:rawbridge}
\end{align}
This follows by the ordered product rule. It is used only when its old smooth
fields and displayed distributional products exist on the common realization.
Already singular higher words require their original tested-kernel formulation;
no product of coincident boundary delta distributions is inserted.

For $H$, changing the score requires
\begin{equation}
r_i=r_{H,i}+e_i,\qquad e_i=t_{H,i}-t_i.
\label{eq:effectivecorrection}
\end{equation}
Thus the covariance in \eqref{eq:boundarybridge} contains the four original
$dd,d\ell,\ell d,\ell\ell$ rows from \eqref{eq:Casimir}, \emph{and}
\begin{equation}
\Cov(r_{H,1},e_2)+\Cov(e_1,r_{H,2})+\Cov(e_1,e_2).
\label{eq:extracovrows}
\end{equation}
The constant Casimir terms remain in raw first-source masses and in separate
boundary means before centering. The field-dependent action row cannot be
removed by centering. Applying a bound to the algebraically complete
$O_H=t_H+r_H$ can avoid losses from separately bounding its artificial pieces,
but does not prove a smaller bound for every old row. Equations
\eqref{eq:changeofsplit}--\eqref{eq:extracovrows} record the exact conversion.

\section{A quantitative comparison of the actual connected native contact}
\label{sec:contactcomparison}

On a genuine fixed-domain chart, let $s_\theta=\partial_\theta\log\rho_\theta$
be the score of its complete actual body. Exact differentiation gives
\begin{equation}
\partial_\theta\log Z_\theta=K_\theta s_\theta,
\qquad
\partial_\theta K_\theta F_\theta
=K_\theta\dot F_\theta+\Cov_{K_\theta}(F_\theta,s_\theta).
\label{eq:parameterresponse}
\end{equation}
In particular, retaining the actual inherited $A_{12}$,
\begin{equation}
\boxed{\dot c_{12}=K\dot A_{12}
+\Cov(\dot O_1,O_2)+\Cov(O_1,\dot O_2)
+K(V_{12}\cdot\nabla s_\theta).}
\label{eq:c12derivative}
\end{equation}
The last term is
$\Cov(A_{12}+(O_1-m_1)(O_2-m_2),s_\theta)$ by the weak Poisson equation.
It retains the full third cumulant and the inherited-contact score response.
Consequently
\begin{align}
|\dot c_{12}|\le{}&K|\dot A_{12}|
+C_P\{\|\nabla\dot O_1\|_2\|\nabla O_2\|_2
+\|\nabla O_1\|_2\|\nabla\dot O_2\|_2\}\notag\\
&+\{C_P\|\nabla A_{12}\|_2
+2c_4C_P^{3/2}\|\nabla O_1\|_4\|\nabla O_2\|_4\}
\|\nabla s_\theta\|_2.
\label{eq:generalresponsebound}
\end{align}
If a selector moves, let $\dot\mu_\theta^\partial$ be its actual raw boundary
or constraint distribution and
$\mathcal S_\theta(F)=Z_\theta^{-1}\langle\dot\mu_\theta^\partial,F\rangle$.
The additional term in \eqref{eq:c12derivative} is exactly
\begin{equation}
\mathcal S_\theta(A_{12}+t_1^Pt_2^P)
-K_\theta(A_{12}+t_1^Pt_2^P)\mathcal S_\theta(1).
\label{eq:movingboundary}
\end{equation}
There is no logarithm of a sharp indicator. Derivatives of the independent scalar
in \eqref{eq:absolutescalar} are still required for raw responses.

Now specialize to the \emph{verified same-target, same-forest, same-selector}
top-scale interpolation of~\cite[\S4]{ND}:
\begin{equation}
K_\theta=\frac{e^{-\theta w}}{K_0e^{-\theta w}}K_0,
\qquad0\le\theta\le1,\qquad s_\theta=-w,
\quad C_P\le\kappa^{-1},\quad
|\nabla w|\le\alpha_w+2\delta_w|B|.
\label{eq:actualtilt}
\end{equation}
The actual fine map and the physical functions $O_i$ remain fixed on this branch.
Let $n_B$ be the actual real native coordinate count and
$L_*=\sup_{0\le\theta\le1}|\nabla[-\log\rho_\theta](0)|$ its actual drift bound.
This verified interpolation also has
$\nabla^2[-\log\rho_\theta]\ge\kappa I$ on the original convex cutoff,
whose outward normal satisfies $B\cdot n\ge0$.
Using this stronger Hessian and boundary input, the moments of~\cite[\S3.1]{ND} are
\begin{align}
M_2^*&=2n_B/\kappa+L_*^2/\kappa^2,\notag\\
M_4^*&=\left(\frac2\kappa\right)^2
\left(n_B+\frac{L_*^2}{2\kappa}\right)
\left(n_B+2+\frac{L_*^2}{2\kappa}\right).
\label{eq:nativemoments}
\end{align}
They bound $K_\theta|B|^2$ and $K_\theta|B|^4$ without a cutoff supremum;
these moment statements do not follow from a Poincar\'e bound alone.
For the fully admitted padded interior $H$ stencils, its proved native response is
\begin{equation}
|\nabla_BO_i|\le A_i(\epsilon_{\rm bg}+h|B|),\qquad
A_i=C_{M,\rm col}|\Gamma_i|\|f_i\|_\infty
\gamma_i^{-2}V_{D,i}h.
\label{eq:actualHgradient}
\end{equation}
Here $h$ is native and $\gamma_i$ is the saved source coupling. Precisely,
$V_{D,i}=a_{{\rm cur},i}^4N_{D,i}$, where $a_{{\rm cur},i}$ is the actual
source/reference-word mesh in the current units of the admitted~\cite[\S3]{ND}
record and $N_{D,i}$ is its actual stencil count. No conversion to a different
global physical-volume convention is implicit. The saved factor $\Gamma_i$ is internally
constant; otherwise its full product-rule derivative must be added. Distinct old
meshes require their actual pullback-response verification before this estimate
can be used.

\begin{theorem}[Actual native two-source contact comparison]
\label{thm:contactvariation}
Under precisely \eqref{eq:actualtilt}--\eqref{eq:actualHgradient}, with unchanged
physical $H$ insertions and bare $A_{12}=0$,
\begin{equation}
\boxed{\begin{aligned}
&|\Cov_{K_1}(O_1,O_2)-\Cov_{K_0}(O_1,O_2)|\\
&\quad\le2c_4\kappa^{-3/2}A_1A_2
\bigl(\epsilon_{\rm bg}+h(M_4^*)^{1/4}\bigr)^2
\bigl(\alpha_w+2\delta_w\sqrt{M_2^*}\bigr).
\end{aligned}}
\label{eq:contactvariation}
\end{equation}
For an unchanged inherited $A_{12}$, the comparison of the actual $c_{12}$
has the additional term
\begin{equation}
\kappa^{-1}(\alpha_w+2\delta_w\sqrt{M_2^*})
\sup_{0\le\theta\le1}\|\nabla_BA_{12}\|_{L^2(K_\theta)}.
\label{eq:inheritedcontact}
\end{equation}
\end{theorem}
\begin{proof}
Minkowski and \eqref{eq:nativemoments} give
$\|\nabla O_i\|_4\le A_i(\epsilon_{\rm bg}+h(M_4^*)^{1/4})$
and $\|\nabla w\|_2\le\alpha_w+2\delta_w\sqrt{M_2^*}$.
Use \eqref{eq:c12derivative} and \eqref{eq:V12bound}, then integrate in $\theta$.
An independent proof avoids Poisson regularity: differentiation gives
$\partial_\theta\Cov(O_1,O_2)
=-\Cov((O_1-m_1)(O_2-m_2),w)$.
Poincar\'e's covariance bound, H\"older, and \eqref{eq:L4Poincare} give the
same constant $2c_4C_P^{3/2}$ directly.
\end{proof}

When $\gamma_1=\gamma_2=h$, the source part of the bound is proportional to
\[
|\Gamma_1\Gamma_2|\|f_1\|_\infty\|f_2\|_\infty V_{D,1}V_{D,2}
\bigl(\epsilon_{\rm bg}/h+(M_4^*)^{1/4}\bigr)^2
(\alpha_w+2\delta_w\sqrt{M_2^*}).
\]
The estimate concerns the complete physical covariance. The old action and
effective-score covariance rows in \eqref{eq:extracovrows} have not individually
lost their inverse mesh factors. Moving insertions, moving contacts, and moving
selectors require the explicit extra terms in
\eqref{eq:generalresponsebound}--\eqref{eq:movingboundary}. This theorem gives
neither all-history MT36 nor its required age and source-grade gains.

\section{A sharper exact Wilson pressure comparison}
\label{sec:pressure}

Make the plaquette convention explicit in this section:
\begin{equation}
s(U)=c_{\rm rep}\left(1-\frac1N\Re\Tr U\right),\qquad c_{\rm rep}>0.
\label{eq:actionconvention}
\end{equation}
$c_{\rm rep}=N$ is \eqref{eq:Wilson}; $c_{\rm rep}=1$ is the SV2 convention.
In either case $\beta$ is its original coefficient and is unchanged in every
volume comparison. For a rectangular four-torus $R$ put
\begin{equation}
Z_{a,R}(\beta)=\int e^{-\beta\sum_{p\subset R}s(U_p)}\prod_b\dd U_b,
\qquad p_{a,R}=\Vol(R)^{-1}\log Z_{a,R}.
\label{eq:pressuredef}
\end{equation}
All action constants and normalized Haar factors are included. Write
$T_*=(\R/t_*\mathbb Z)^4$, $n=t_*/a$ and an integer cover factor $q$.

\subsection{The exact positive transfer operator and its gauge projection}

For a transverse three-torus $\Sigma$ of physical volume $A$, its original
numbers of sites and positive spatial links are
\begin{equation}
V_\Sigma=A/a^3,\qquad E_\Sigma=3A/a^3.
\label{eq:slicecount}
\end{equation}
On $X=G^{E_\Sigma}$ with normalized Haar, set
\begin{align}
b(U)&=e^{-\beta S_{\rm sp}(U)/2},\qquad k_\beta(g)=e^{-\beta s(g)},\notag\\
C_\beta(U,V)&=\prod_{e\in E_\Sigma}k_\beta(U_eV_e^{-1}),\qquad
T=M_bC_\beta M_b.
\label{eq:transfer}
\end{align}
In particular $k_\beta(1)=1$; its factor $e^{-\beta c_{\rm rep}}$ per temporal
plaquette has not been divided out. The pure-gauge transfer construction and
its character-positivity mechanism agree with the original construction of
L\"uscher~\cite{luscher}; the periodic projected identity is derived below.

\begin{lemma}[Positivity, trace, and physical top eigenvalue]
\label{lem:transfer}
$T$ is positive, self-adjoint and trace class. Let $P$ be the normalized Haar
average over site gauge transformations on $\Sigma$. Then $TP=PT$, and
\begin{align}
\Tr T&=\int_Xb(U)^2\dd U,\notag\\
Z_{a,(ra)\times\Sigma}&=\Tr(T^rP)
=\Tr_{\rm phys}T_{\rm phys}^r,\qquad
T_{\rm phys}=T|_{PL^2(X)},\label{eq:periodictrace}\\
\|T_{\rm phys}\|&=\|T\|=:\lambda>0.
\label{eq:physicaltop}
\end{align}
\end{lemma}
\begin{proof}
The exact absolutely convergent series
\[
k_\beta(g)=e^{-\beta c_{\rm rep}}
\sum_{r,s\ge0}\frac{(\beta c_{\rm rep}/(2N))^{r+s}}{r!s!}
\chi_f(g)^r\chi_{\bar f}(g)^s
\]
has nonnegative irreducible character coefficients $c_\rho$: each product is
the character of a tensor representation. Convolution has eigenvalue
$c_\rho/d_\rho$ on its Peter--Weyl block of multiplicity $d_\rho^2$.
The nonnegative total trace is
$\sum_\rho d_\rho c_\rho=k_\beta(1)=1$.
Tensoring over the actual spatial links and multiplying by bounded $b$ proves
positivity and trace class. The continuous diagonal kernel gives the stated trace.

For a slice gauge field $h$, $(h\cdot U)_{xy}=h_xU_{xy}h_y^{-1}$.
The class function and spatial action are invariant under simultaneous changes
of $U,V$, so $T$ commutes with $P$. Keeping the temporal links between two slices,
the temporal plaquette at $e=(x,y)$ is
$U_eh_yV_e^{-1}h_x^{-1}=U_e(h\cdot V)_e^{-1}$.
Their exact Haar integration gives the layer kernel $TP$. Spatial action halves
combine once per slice. Closing $r$ layers gives $\Tr((TP)^r)=\Tr(T^rP)$.
This keeps the temporal holonomy integral and introduces no gauge-volume scalar.

Finally compact positivity gives a top Rayleigh eigenvector $f$.
The strictly positive kernel implies
$\langle|f|,T|f|\rangle\ge\langle f,Tf\rangle=\lambda\|f\|^2$,
so $|f|$ is also in the top eigenspace. Its nonnegative gauge average $P|f|$
is nonzero, since it has the same positive integral. Commutation then gives a
physical top eigenvector with the same eigenvalue, proving \eqref{eq:physicaltop}.
No uniqueness or weak-coupling assumption is required.
\end{proof}

\subsection{A covering bound that costs only logarithmically in $\beta$}

Fix a bi-invariant group metric. For $0<r\le r_0$ there is a measurable group
partition into at most $C_{\rm cov}r^{-m}$ sets of diameter at most $r$.
One construction is a maximal separated net followed by a partition subordinate
to its ball cover. Smoothness at the action minimum gives
$s(uv^{-1})\le C_sr^2$ for points in one such set. Take
$r=r_0(1+\beta)^{-1/2}$ and form the product partition of $X$.
For $U,V$ in one product cell $D$,
\begin{equation}
C_\beta(U,V)\ge e^{-C_sr_0^2E_\Sigma}.
\label{eq:cellkernel}
\end{equation}
Some cell has
$m_D:=\int_Db^2\ge\Tr T/(C_{\rm cov}r^{-m})^{E_\Sigma}$.
The exact Rayleigh trial $f_D=b\ind_D/\sqrt{m_D}$ gives
\begin{align}
\lambda&\ge\frac1{m_D}\int_{D\times D}
b(U)^2b(V)^2C_\beta(U,V)\dd U\dd V\notag\\
&\ge e^{-C_sr_0^2E_\Sigma}m_D.
\label{eq:Rayleigh}
\end{align}
Consequently, for a fixed $C_{G,c_{\rm rep}}$,
\begin{equation}
\boxed{\log\frac{\Tr T}{\|T\|}\le E_\Sigma B_G(\beta),\qquad
B_G(\beta)=C_{G,c_{\rm rep}}+\frac m2\log(1+\beta).}
\label{eq:coverentropy}
\end{equation}
This is valid for the actual spatial action. The unprojected covering inequality
also works for any nonzero bounded nonnegative factor in place of $b$; gauge
invariance remains necessary for the projected identities. That breadth explains
why it does not yet exploit spatial dynamics enough to remove the fine-link
surface count.

Using positivity and \eqref{eq:physicaltop}, for every integer $r\ge1$,
\begin{equation}
1\le\Tr_{\rm phys}(T_{\rm phys}/\lambda)^r
\le\frac{\Tr T}{\lambda},\qquad
\log\Tr_{\rm phys}\frac{T_{\rm phys}^r}{\|T_{\rm phys}^r\|}
\le\frac{3A}{a^3}B_G(\beta).
\label{eq:traceentropy}
\end{equation}
The upper bound estimates the physical trace by the unprojected trace; this loss
is explicit, while the exact gauge projection is retained in the partition function.

\begin{theorem}[Logarithmic seam pressure bound]
\label{thm:pressure}
For the original Wilson action at the same $\beta\ge0$, and every compatible
integer $q\ge1$,
\begin{equation}
\boxed{0\le p_{a,T_*}-p_{a,qT_*}
\le\frac{12}{t_*a^3}
\left[C_{G,c_{\rm rep}}+\frac m2\log(1+\beta)\right].}
\label{eq:pressurebound}
\end{equation}
\end{theorem}
\begin{proof}
For fixed $\Sigma$ set $K=T_{\rm phys}^{n}$, $n=t_*/a$, and
$\widehat K=K/\|K\|$. The same-transverse identity \eqref{eq:periodictrace} gives
\begin{equation}
p_{a,t_*\times\Sigma}-p_{a,qt_*\times\Sigma}
=\frac1{t_*A}\left[
\log\Tr\widehat K-\frac1q\log\Tr\widehat K^q\right].
\label{eq:rectangularpressure}
\end{equation}
The top eigenvalue cancels exactly within this pair. Its bracket lies between
zero and $\log\Tr\widehat K$, hence the difference is at most
$3B_G(\beta)/(t_*a^3)$ by \eqref{eq:traceentropy}.
Enlarge the four sides successively. The $j$th transverse volume is
$A_j=q^{j-1}t_*^3$, with exactly $3A_j/a^3$ spatial links, and the enlarged
side starts at $t_*$. Summing the four bounds proves \eqref{eq:pressurebound}.
The four transverse top eigenvalues need not agree; each cancels only in its
own exact pair. No volume-dependent vacuum normalization is introduced.
\end{proof}

Along $\beta_a=O(\log(1/a))$, the estimate is
$O(a^{-3}\log\log(1/a))$. It improves the earlier $O(\beta_a a^{-3})$ bound
but is not ultraviolet uniform. Given the already stated original stability
$|\log Z_{a,R}-V_{a,R}|\le C_{\rm st}\Vol(R)$, the corresponding vacuum
comparison has this bound plus $2C_{\rm st}$. The proof compares complete
partition functions, including all original histories and scalar weights;
it is not a termwise conditional history estimate.

\subsection{A proved energy--entropy cancellation for the actual seam}

Use the original common-variable gluing interpolation
$\Delta S=S_{\rm big}-S_{\rm per}$ and
$\dd\mu_t=Z_t^{-1}e^{-\beta(S_{\rm per}+t\Delta S)}\dd\lambda$.
Let $M=q^4$ and
\begin{equation}
G_{a,q}=M\log Z_{a,T_*}-\log Z_{a,qT_*}
=Mt_*^4(p_{a,T_*}-p_{a,qT_*}).
\label{eq:seamG}
\end{equation}
For $n\ge2,q\ge2$ the original seam count and product-law expectation are
\begin{equation}
N_{\rm seam}=M(12n^3-6n^2),\qquad
\E_0\Delta S=N_{\rm seam}(c_{\rm rep}-\bar s_{a,*}).
\label{eq:seamcount}
\end{equation}
Independent block gauge symmetry makes every new open-loop fundamental trace
average zero, giving $c_{\rm rep}$, not automatically one. The relative entropy
therefore satisfies the exact identity
\begin{equation}
D(\mu_0\Vert\mu_1)
=\beta N_{\rm seam}(c_{\rm rep}-\bar s_{a,*})-G_{a,q}.
\label{eq:seamentropy}
\end{equation}
Theorem~\ref{thm:pressure} gives
$0\le G_{a,q}\le12Mn^3B_G\le\frac43B_GN_{\rm seam}$. Hence
\begin{equation}
\boxed{\begin{aligned}
N_{\rm seam}\left[\beta(c_{\rm rep}-\bar s_{a,*})-\tfrac43B_G(\beta)\right]
&\le D(\mu_0\Vert\mu_1)\\
&\le\beta N_{\rm seam}(c_{\rm rep}-\bar s_{a,*}).
\end{aligned}}
\label{eq:seamcancellation}
\end{equation}
Thus the artificial linear-$\beta$ gluing energy is matched by actual relative
entropy up to a logarithmic surface remainder. This needs no weak-plaquette
asymptotic and does not assume intermediate measures have endpoint gauge symmetry.

The exact remaining cancellation can also be written
\begin{equation}
G_{a,q}=\beta N_{\rm seam}(c_{\rm rep}-\bar s_{a,*})
-\beta^2\int_0^1(1-t)\Var_{\mu_t}(\Delta S)\dd t.
\label{eq:variancecancellation}
\end{equation}
This is Taylor's theorem for $\log Z_t$, whose second derivative is
$\beta^2\Var_{\mu_t}(\Delta S)$. A uniform upper bound on $G$ requires a
matching \emph{lower} bound on this variance integral to accuracy
$O(Mt_*^4/\beta^2)$. An upper variance estimate alone has the wrong sign for
that purpose.
\section{An exact two-copy target for the remaining pressure gap}
\label{sec:pressuretarget}

The preceding argument identifies a concrete remaining input without choosing
any new vacuum normalization. Let $\Sigma$ range over the three-dimensional
rectangular tori used in the four successive enlargements: zero, one, two, or
three sides equal to $qt_*$, with all other sides $t_*$. The cutoff $a$ is
compatible with every side. All constants below must be uniform in $a$, in
integer $q\ge2$, and in this actual transverse family.

For $K=T_{\rm phys}^{t_*/a}$ put
$\widehat K=K/\|K\|$ and $R_{a,\Sigma}=\Tr\widehat K$. Its eigenvalues
$a_j$ lie in $[0,1]$ and include one. Hence
$1\le\sum_j a_j^q\le\sum_j a_j=R_{a,\Sigma}$, and
\begin{equation}
\left(1-\frac1q\right)\log R_{a,\Sigma}
\le\log R_{a,\Sigma}-\frac1q\log\Tr\widehat K^q
\le\log R_{a,\Sigma}.
\label{eq:traceequivalence}
\end{equation}
The remaining improvement needed in \eqref{eq:traceentropy} is therefore
\begin{equation}
\boxed{\log R_{a,\Sigma}\le C_{t_*}\Vol_3(\Sigma).}
\label{eq:physicaltrace}
\end{equation}
This is a hypothesis still to be proved for the original interacting operator.

For clarity, it is quantitatively equivalent to the ultraviolet-uniform
\emph{pressure part} of the desired periodic comparison on this family.
If \eqref{eq:physicaltrace} holds, \eqref{eq:rectangularpressure} and four
enlargements give $p_{a,T_*}-p_{a,qT_*}\le4C_{t_*}/t_*$.
Conversely suppose this isotropic pressure difference is at most $C_p$.
Every intermediate rectangular difference is nonnegative, so each is at most
$C_p$. Its own \eqref{eq:traceequivalence} gives
\[
\log R_{a,\Sigma}\le\frac q{q-1}t_*A C_p\le2t_*A C_p.
\]
The reference transverse torus is already covered by the first step at $q=2$.
Any additional comparison with the original vacuum functional still uses the
stated absolute stability estimate, with its unchanged constant.

Writing the actual physical eigenvalues of $T_{\rm phys}$ as $\lambda_j$ and
its top eigenvalue as $\lambda_0$, define
\begin{equation}
E_{a,\Sigma,j}=-a^{-1}\log(\lambda_j/\lambda_0),\qquad
\log R_{a,\Sigma}=\log\sum_j e^{-t_*E_{a,\Sigma,j}}.
\label{eq:excitationcount}
\end{equation}
A zero transfer eigenvalue contributes zero to the sum. The missing bound
controls the number and energies of physical excitations at fixed physical time.
A positive mass gap alone would not give this estimate. For example, an
abstract finite-dimensional operator with one eigenvalue one and $R_a-1$
eigenvalues $e^{-a\Delta}$ has gap $\Delta>0$, but its normalized time-$t_*$
trace is $1+(R_a-1)e^{-t_*\Delta}$. If $\log R_a$ is of order $A/a^3$,
\eqref{eq:physicaltrace} fails. This diagnostic is not a Wilson counterexample.

\subsection{The equivalent original two-copy overlap}

The top eigenvalue can be eliminated from the target. Work on the original
product-Haar boundary space $X$, and set
$K=(TP)^{t_*/a}$, extended by zero on the orthogonal complement of the
physical subspace. Its kernel includes the original temporal integrations and
gauge averages. Define
\begin{align}
\mathcal P_{a,\Sigma}
&=\frac{\Tr K^2}{(\Tr K)^2}
=\frac{Z_{a,2t_*\times\Sigma}}{Z_{a,t_*\times\Sigma}^2},
\notag\\
\mathcal E_{2,a,\Sigma}&=-\log\mathcal P_{a,\Sigma}.
\label{eq:purity}
\end{align}
These are exact original partition functions. From
$1\le\Tr\widehat K^2\le R_{a,\Sigma}$ one obtains
\begin{equation}
\log R_{a,\Sigma}\le\mathcal E_{2,a,\Sigma}
\le2\log R_{a,\Sigma},\qquad
p_{a,t_*\times\Sigma}-p_{a,2t_*\times\Sigma}
=\frac{\mathcal E_{2,a,\Sigma}}{2t_*A}.
\label{eq:purityequivalence}
\end{equation}
Consequently the following explicit inequality is equivalent, up to fixed
constants, to \eqref{eq:physicaltrace}:
\begin{equation}
\boxed{
\int_{X\times X}K(U,V)K(V,U)\dd U\dd V
\ge e^{-C_{t_*}A}
\left(\int_XK(U,U)\dd U\right)^2.}
\label{eq:overlapneeded}
\end{equation}
The real symmetric kernel makes the left integrand $K(U,V)^2$.
A common scalar multiplying this same kernel cancels from the ratio.
Independent normalizations of the one-slab and two-slab partition functions
would change it and are not allowed.

For an actual nonnegative partition of the original slab integral into
histories, $K=\sum_{\mathfrak h}K_{\mathfrak h}$ with the original raw kernels,
the left side of \eqref{eq:overlapneeded} is exactly
\begin{equation}
\sum_{\mathfrak h,\mathfrak h'}
\int K_{\mathfrak h}(U,V)K_{\mathfrak h'}(V,U)\dd U\dd V.
\label{eq:historyoverlap}
\end{equation}
Every original scalar and history weight remains in this expression. A signed
connected polymer expansion is not automatically a nonnegative partition.
Restricting to matching histories needs a quantitative bound on the mass lost.
Any proposed coarse change of variables must preserve the original seam
identifications and coarea factors in both integrals.

The proved estimate is only
\begin{equation}
\mathcal E_{2,a,\Sigma}\le
\frac{6A}{a^3}\left[C_{G,c_{\rm rep}}+\frac m2\log(1+\beta)\right].
\label{eq:purityproved}
\end{equation}
Replacing its fine-link count by physical transverse volume is the outstanding
interacting estimate. Neither a free determinant nor a separately bounded
$O(g_a^p)$ error per fine site at any fixed power $p$ supplies it: division
by physical volume gives $O(a^{-4}g_a^p)$, still divergent on the logarithmic
trajectory. One needs a bound on the actual \emph{difference} of vacuum
corrections or the exact original overlap \eqref{eq:overlapneeded}.

\section{Verified scope, remaining problems, and checks}
\label{sec:status}

\subsection{What this note establishes}

The local Wilson moment proof reuses the original matching partition powers,
checks their absolute gauge count, and extends the cell argument to every even
torus with its multiplicity four. It yields the explicit scalar
$\mathcal B_{n,\beta}$, improves the original compensated $H$ contact bound,
and strengthens the fixed-patch bad-bath probability. These are statements
about the original Wilson law, not about arbitrary conditioned histories.

The conditional prescription \eqref{eq:newsource} is an explicit new
finite-cutoff hierarchy. Its local contacts have proved finite-order absolute
masses, including the nonvanishing order $a^{-8}$ two-source subtraction.
The correlated-pair test proves why that one pass does not establish a tight
remaining field. The separate Gaussian extension has an explicit contact
order, exact logarithmic residue, and a proved ordinary-covariance obstruction
at that order. These conclusions do not substitute for the interacting hierarchy.

The one-step chart proves admission for the specified literal local plaquette
selector and original read set, plus exact current/readout lifts and coarea
scores. The Poisson identities use the complete actual effective density on
an identified positive fiber. On the previously verified native interpolation,
the connected physical two-source comparison \eqref{eq:contactvariation} is
a concrete new estimate, with actual source units and native coordinate counts.
Its exact conversion to the old representation retains the Casimir rows,
effective-score corrections, sharp-boundary partners, and absolute normalizers.

The exact Wilson transfer argument improves the pressure cost from linear
to logarithmic in the bare coupling at fixed fine surface area and proves
the corresponding seam-energy/entropy cancellation. It then reduces the
remaining periodic pressure requirement to the explicit original
two-copy inequality \eqref{eq:overlapneeded}. That inequality itself is unproved.

\subsection{The unresolved estimates}

\begin{enumerate}
\item \textbf{Interacting contact completion.}
Prove the needed single-operator mixing and the interacting off-diagonal
estimate \eqref{eq:interactingneeded}, or another estimate sufficient for a
specified replacement prescription. At every source order, construct compatible
extensions at all partial diagonals and verify the required positivity and
continuum identities. A finite-cutoff identity, an averaged contact-mass bound,
or a two-point Gaussian extension does not provide these facts. The
field-dependent Casimir action row must remain in the original-source calculation.

\item \textbf{Actual general mixed-selector/readout admission.}
Identify the original $R$-I/$R$-II selectors, their parameter ranges
$(1.64)$--$(1.69)$ and $(1.75)$, and the three joins on the same actual
realization used by AE8/LM47. Verify the old readout and moving-reference
response there, retaining the nonabelian relative-target correction and the
sharp joint constraint distributions. The one-step local chart here is not
that general identification. No lemma from the $R$-I paper beyond the cited
equations is presumed.

\item \textbf{All-history physical-source transfer.}
For each required positive body, prove its actual domain and spectral estimate,
or use another justified transport. Bound the actual inherited $O_i,A_I$,
their native and retained-reference derivatives, the readout differential,
and the original boundary distributions in the required norms. Keep every
external scalar and the source-grade and history-age gains required by MT36.
The native estimate proved here covers its stated fixed-target interpolation;
the abstract inverse identity does not verify those hypotheses in other histories.

\item \textbf{Ultraviolet-uniform pressure comparison.}
Establish \eqref{eq:overlapneeded}, equivalently \eqref{eq:physicaltrace},
on the actual transverse family, or prove the needed all-scale vacuum
cancellation with the original absolute normalizers. The logarithmic surface
estimate \eqref{eq:pressurebound} still diverges as $a\to0$.

\item \textbf{The final continuum construction.}
Assemble a common noncollapsing continuum hierarchy with the required Euclidean
and Osterwalder--Schrader properties, compatible physical source identifications,
and the intended coupling trajectory. None of the four preceding gaps may be
replaced by a conditional assertion in that construction. This note
does not prove Yang--Mills existence or a mass gap.
\end{enumerate}

\subsection{Checks and their limits}

An exact finite calculation uses only
integer and rational arithmetic. It constructs the full $12\times12$
directional-current covariance matrix from the fourth derivative of $|x|^{-2}$, verifies the
total-current cancellations modulo $|\omega|^2=1$, checks
\eqref{eq:angular}, and checks all 165 degree-eight monomial coefficients in
\eqref{eq:P8explicit}. No sampled integration or floating-point inference is used.

\begin{thebibliography}{9}
\bibitem{W}
Wilson-source estimates, volumes 1--10 (\texttt{lattice\_theory/source\_estimates} in this collection), cited by source page. Matching partition
powers and chessboard moments: \S10, pp.~21--23.

\bibitem{ND}
\emph{Nonlinear dispersion and unequal spatial weights}, \texttt{lattice\_theory/block\_and\_transfer/conditional\_estimates/nonlinear\_dispersion\_and\_unequal\_spatial\_weights.tex} in this collection. Used here for the complete nonlinear
one-link non-tightness argument, the admitted native Poincar\'e/moment and
source-response records, and the original boundary/normalizer bookkeeping.

\bibitem{biskup}
M. Biskup, \emph{Reflection Positivity and Phase Transitions in Lattice Spin
Models}, lecture notes. Closed-cell setup and Theorem 5.8,
pp.\ 43--46; equations (5.34)--(5.35).
\url{https://www.math.ucla.edu/~biskup/PDFs/papers/Biskup-Prague-notes.pdf}.

\bibitem{bf}
R. Brunetti and K. Fredenhagen,
\emph{Microlocal Analysis and Interacting Quantum Field Theories:
Renormalization on Physical Backgrounds},
Commun.\ Math.\ Phys.\ \textbf{208} (2000), 623--661.
Scaling-degree extension results: \S5, especially Theorem 5.3.
\url{https://arxiv.org/abs/math-ph/9903028}.

\bibitem{luscher}
M. L\"uscher, \emph{Construction of a Selfadjoint, Strictly Positive Transfer
Matrix for Euclidean Lattice Gauge Theories}, DESY 76/54 (1976).
Pure-gauge factors (17), (20), (21), Proposition 1, and character argument
(22)--(23).
\url{https://bib-pubdb1.desy.de/record/396349/files/7611148.pdf}.
\end{thebibliography}
\end{document}
